% =====================================================================
\section{Conclusiones}
\label{sec:conclusiones}

\subsection*{Conclusión 1 --- Unidad I: el contexto determinó el alcance más que la técnica}

La planificación de la elicitación fijó como objetivo cubrir a los seis perfiles de usuario, y esa
decisión ---tomada antes de conocer el dominio--- resultó ser la que más restricciones generó. El
cuestionario ampliado \id{EV-12}, con $n=62$, reveló que el 11,3\,\% del personal no dispone de
teléfono inteligente y que el 74,2\,\% no tiene señal permanente en el lote. Ambos datos
contradijeron suposiciones que el equipo había dado por válidas y forzaron dos requisitos que no
estaban previstos: el registro delegado \id{RF-35} y la operación sin conexión \id{RD-02}.

La conclusión operativa es que la cobertura de la elicitación no es una cuestión de exhaustividad
metodológica sino de alcance del sistema: si el instrumento no hubiera llegado al personal de campo,
el ERS habría especificado un sistema que el 11,3\,\% de sus usuarios no podría utilizar, y el
defecto no habría aparecido hasta el despliegue.

\subsection*{Conclusión 2 --- Unidad II: la evidencia tardía es la más cara}

La evidencia \id{EV-13} ---la hoja de liquidación semanal--- se obtuvo después de haber cerrado la
especificación, y obligó a incorporar cuatro requisitos funcionales nuevos sobre un proceso que la
organización ejecuta todas las semanas y que consume una jornada administrativa completa. Ninguna de
las doce evidencias anteriores lo había identificado.

El costo de esa incorporación tardía es medible y se distribuyó por todo el documento: obligó a
modificar recuentos en cuatro archivos, produjo la contradicción entre 21 y 19 requisitos
\emph{Must} que la inspección detectó como defecto crítico, y dejó dos requisitos \emph{Must} sin
historia de usuario que hubo que tramitar ante el Change Control Board. Un requisito que llega tarde
no cuesta lo que cuesta especificarlo: cuesta lo que cuesta reconciliar el documento entero con él.

\subsection*{Conclusión 3 --- Unidad III: dos taxonomías del mismo sistema producen defectos}

Once de los treinta y tres defectos detectados en la inspección son de consistencia, y ninguno es de
redacción. El patrón tiene una causa identificable: el documento mantenía la parte documental
---bloques funcionales, fichas de requisito--- y la parte de priorización ---MoSCoW, WSJF, matriz---
como dos representaciones del mismo conjunto de requisitos, sin mecanismo que las reconciliara.

La Unidad V confirmó el diagnóstico desde otro ángulo. Al incorporar los diagramas de flujo de datos,
la decisión deliberada fue hacer corresponder uno a uno los siete procesos con los siete bloques
funcionales existentes, en lugar de descomponer el sistema por criterio propio del análisis de flujo.
El modelo resultante es posiblemente menos fiel que el que habría producido una descomposición
independiente, y se aceptó esa pérdida precisamente para no crear una tercera taxonomía paralela. El
acoplamiento medio de 2,20 de la métrica M5 es la consecuencia medible de esa decisión.

\subsection*{Conclusión 4 --- Unidad IV: la inspección es barata y el alcance que se le fija es lo caro}

La inspección formal detectó treinta y tres defectos a un costo de 0,44 horas-persona por defecto, lo
que sostiene empíricamente el argumento de Fagan sobre el costo relativo de la detección temprana
\cite{fagan1976}. Pero la re-inspección de la Unidad V encontró siete defectos residuales, y su
distribución es más informativa que su número: cuatro se concentran en los apéndices y en los
conteos globales, que el alcance de la inspección había excluido de forma expresa.

La exclusión era razonable cuando se decidió ---los apéndices son material de referencia--- y dejó de
serlo en cuanto tres solicitudes de cambio modificaron los recuentos que allí residen. Nadie revisó
la decisión de alcance después de aprobar las RFC. La conclusión no es que la inspección fallara,
sino que un alcance de inspección es una decisión con fecha de caducidad: caduca cuando cambia el
documento sobre el que se fijó.

\subsection*{Conclusión 5 --- Unidad V: especificar equidad exige preguntar entre quiénes}

El ERS de la Entrega 2A especificaba exactitud, latencia y explicabilidad para sus componentes
automáticos, y no especificaba equidad. La ausencia no fue un descuido de redacción: fue un punto
ciego del modelo de calidad empleado, que no contemplaba la propiedad, y solo se hizo visible al
formular explícitamente la pregunta de entre qué grupos podía diferir el desempeño del sistema.

La respuesta, en un cultivo de palma, no es demográfica. Los cuatro requisitos de equidad
especificados comparan variedades sembradas, gamas de dispositivo y cuadrillas de trabajo. El de
mayor consecuencia es \id{RNF-IA-10}: como \id{RF-24} traza la penalización de la extractora hasta la
cuadrilla responsable, un modelo que estimara peor las cargas de una cuadrilla concreta produciría
recortes de remuneración sistemáticamente injustos sobre personas identificables. Es el único
indicador del plan de monitoreo cuya acción al superarse es suspender el uso del modelo y no
reentrenarlo.

La conclusión que integra las cinco unidades es esta: el proyecto midió lo que sabía nombrar. Los
defectos de consistencia aparecieron cuando se contaron las cifras; los casos de uso faltantes,
cuando se dividió especificados entre identificados; y la ausencia de requisitos de equidad, cuando
se preguntó entre qué grupos. En los tres casos el problema existía antes de la medición y llevaba
meses en el documento sin que ninguna lectura lo revelara. La ingeniería de requisitos no consiste
solo en escribir bien lo que se sabe, sino en construir los instrumentos que hacen visible lo que
todavía no se ha nombrado.
